\documentclass[11pt]{amsart}
\usepackage{amsmath,amssymb,amsthm}
\usepackage{booktabs}
\usepackage[colorlinks=true,allcolors=blue,linktoc=page]{hyperref}
\usepackage[capitalize]{cleveref}   % without the option nameinlink, only the number is a link

\theoremstyle{plain}
\newtheorem{theorem}{Theorem}[section]
\newtheorem{proposition}[theorem]{Proposition}
\newtheorem{lemma}[theorem]{Lemma}
\newtheorem{corollary}[theorem]{Corollary}
\theoremstyle{definition}
\newtheorem{definition}[theorem]{Definition}
\newtheorem{remark}[theorem]{Remark}
\newtheorem{example}[theorem]{Example}
\newtheorem{question}[theorem]{Question}

\title[Three layers of cyclotomic $\mathbb{Z}_p$-extensions with class number one]{Three layers of cyclotomic $\mathbb{Z}_p$-extensions of $\mathbb{Q}$ with class number one}
% \author{}   % to be filled in
\date{November 26, 2026 (third draft)}

\emergencystretch=1.5em

\begin{document}

\begin{abstract}
We prove that the following three fields have class number $1$: the maximal real subfield of the cyclotomic field of conductor $243$, the second layer of the cyclotomic $\mathbb{Z}_7$-extension of $\mathbb{Q}$, and the first layer of the cyclotomic $\mathbb{Z}_{29}$-extension of $\mathbb{Q}$.
The proofs do not assume the generalized Riemann hypothesis.
They combine an upper bound for class numbers due to Miller with a list of elements of small height whose norms are factored.
From this list we deduce that the prime ideals above $16049777$ prime numbers are principal in the first field, above $29835$ prime numbers in the second field, and above $6549261$ prime numbers in the third field.
We also prove that the first layers of the cyclotomic $\mathbb{Z}_p$-extensions of $\mathbb{Q}$ for $p = 37$, $41$, and $43$ have class number $1$ under the generalized Riemann hypothesis.
\end{abstract}

\maketitle

\section{Introduction}

For each positive integer $m$, let $k_m$ denote the unique subfield of degree $m$ of the cyclotomic $\hat{\mathbb{Z}}$-extension of $\mathbb{Q}$.
If $p$ is a prime number and $n$ is a positive integer, then $k_{p^n}$ is the $n$-th layer of the cyclotomic $\mathbb{Z}_p$-extension of $\mathbb{Q}$.
This field is denoted by $\mathbb{B}_{p,n}$ in \cite{Mi15b}, and the field $k_{2^n}$ is often denoted by $\mathbb{B}_n$.
Let $h(k_m)$ denote the class number of $k_m$.

Weber \cite{We86} proved that $h(k_{2^n})$ is odd for every $n$ and that $h(k_{2^n}) = 1$ if $n \le 3$.
Whether $h(k_{2^n}) = 1$ holds for every $n$ is known as Weber's class number problem.
Coates \cite{Co12} asked whether $h(k_{p^n}) = 1$ holds for every prime number $p$ and every positive integer $n$.
At present, the equality $h(k_{p^n}) = 1$ is proved without any hypothesis for
\[
p^n \in \{2, 2^2, 2^3, 2^4, 2^5, 2^6,\ 3, 3^2, 3^3,\ 5, 5^2,\ 7,\ 11,\ 13,\ 17,\ 19\},
\]
and it is proved under the generalized Riemann hypothesis for
\[
p^n \in \{2^7,\ 3^4,\ 23,\ 29,\ 31\};
\]
these results are due to Weber, Cohn, Bauer, Masley, van der Linden, Coates, Liang, Mih\u{a}ilescu, and Miller; see \cite[Table 1, Theorem 7, and Table 2]{Mi15b} and \cite{Co60, Ba69, Ma78, vdL82, Co12, Mi14, Mi14b}.
For a fixed prime number $\ell$, the question whether $\ell$ divides $h(k_{p^n})$ for some $n$ has been studied by Horie and Horie \cite{HH09}, Fukuda and Komatsu \cite{FK11}, Morisawa and Okazaki \cite{MO13}, and Fukuda, Komatsu, and Morisawa \cite{FKM14}, among others.
An unconditional proof of $h(k_{23}) = 1$ is claimed in \cite{Ch26}; to our knowledge, it has not yet appeared in a refereed publication.

The field $k_{3^n}$ is the maximal real subfield $\mathbb{Q}(e^{2\pi i/3^{n+1}})^+$ of the cyclotomic field of conductor $3^{n+1}$.
In this note we prove the following theorem.

\begin{theorem} \label{thm:main}
\begin{enumerate}
\item The class number of $k_{3^4} = \mathbb{Q}(e^{2\pi i/243})^+$ is $1$.
\item The class number of $k_{7^2}$ is $1$.
\item The class number of $k_{29}$ is $1$.
\end{enumerate}
\end{theorem}

The proof does not assume the generalized Riemann hypothesis.
Under this hypothesis, the equality $h(k_{3^4}) = 1$ was proved by van der Linden \cite[Theorem 2]{vdL82}, and the equality $h(k_{29}) = 1$ was obtained by Miller \cite[Table 2]{Mi15b}.

We also prove a conditional result for three fields for which, to our knowledge, the class number was not known even under the generalized Riemann hypothesis.

\begin{theorem} \label{thm:grh}
Let $p \in \{37, 41, 43\}$.
If every nontrivial zero of the Dedekind zeta function of the Hilbert class field of $k_p$ has real part $1/2$, then the class number of $k_p$ is $1$.
\end{theorem}

The proof of \cref{thm:grh} is given in \cref{sec:grh}.
It uses the conditional form of the inequality of Miller, which requires only few prime numbers.
The difficulty is to prove that the prime ideals above the least split primes are principal, because the norms of the elements of small height are large.
We prove this for several thousand prime numbers at the same time, by linear algebra over the ring $\mathbb{Z}[e^{2\pi i/p}]$; this part of the proof does not depend on any hypothesis.

The proof follows the method of Miller \cite{Mi14, Mi15b}.
Let $K$ denote one of the three fields.
First, an analytic inequality gives an upper bound for $h(K)$, provided that the prime ideals above sufficiently many prime numbers that split completely in $K$ are known to be principal.
Second, a result on the action of the Galois group on the class group shows that $h(K)$ is either $1$ or at least $t$, where $t = 163$ if $K = k_{3^4}$, $t = 197$ if $K = k_{7^2}$, and $t = 59$ if $K = k_{29}$.
Third, we show that the prime ideals above the following prime numbers are principal: every prime number $q < 3 \cdot 10^{10}$ that splits completely in $k_{3^4}$, every prime number $q \le 2.2 \cdot 10^7$ that splits completely in $k_{7^2}$, and every prime number $q < 4 \cdot 10^9$ that splits completely in $k_{29}$.
With these prime ideals, the upper bound in the first step is smaller than $t$.

The third step differs from \cite{Mi15b}, where a generator is given for each prime ideal.
Here we do not compute generators.
We use elements $\alpha$ of the form $c + \sum_i a_i \eta_i$, where the $\eta_i$ are Gaussian periods and $c$ and the $a_i$ are small integers.
If a prime number $q$ divides the norm of $\alpha$ exactly once, and if the prime ideals above the other prime factors of the norm are already known to be principal, then the prime ideals above $q$ are principal.
For each prime number $q$, such an element is found by lattice reduction, and the assertions on $\alpha$ are verified with exact integer arithmetic.
The same principle underlies the algorithms that compute class groups by collecting relations.
Relations of this form are also used in \cite{Ch26} for the field $k_{23}$; there, the elements are found by an exhaustive search over the elements with coefficients in $\{-1, 0, 1\}$.
Since the upper bound for $h(K)$ comes from the analytic inequality, we need neither the unit group nor the regulator of $K$.

The lists of elements are too large to be printed.
They are available, together with the programs that verify every assertion of \cref{sec:data,sec:bound}, at
\begin{center}
\url{https://github.com/uekijun46/Weber_class_number}
\end{center}

\section{The fields} \label{sec:field}

Let $p$ be an odd prime number and let $n$ be a positive integer.
Put $N = p^n$ and $M = p^{n+1}$, and let $\zeta$ denote the complex number $e^{2\pi i/M}$.
The group $(\mathbb{Z}/M\mathbb{Z})^\times$ is cyclic of order $(p-1)N$.
Let $H$ denote its subgroup of order $p - 1$, so that
\[
H = \{a \in (\mathbb{Z}/M\mathbb{Z})^\times \mid a^{p-1} = 1\}.
\]
Let $K$ denote the subfield of $\mathbb{Q}(\zeta)$ fixed by $H$, where $a \in H$ acts by $\zeta \mapsto \zeta^a$.
Then $K = k_{p^n}$.
The field $K$ is totally real since $-1 \in H$, and it is cyclic of degree $N$ over $\mathbb{Q}$.
Its subfields are the fields $k_{p^j}$ with $0 \le j \le n$, where $k_1 = \mathbb{Q}$.

Let $g$ be a generator of $(\mathbb{Z}/M\mathbb{Z})^\times$.
For $i \in \mathbb{Z}/N\mathbb{Z}$, define the Gaussian period
\[
\eta_i = \sum_{a \in H} \zeta^{g^i a} .
\]
This is well defined since $g^N \in H$.
Each $\eta_i$ is an element of the ring of integers $O_K$ of $K$.
Let $\sigma$ denote the restriction to $K$ of the automorphism $\zeta \mapsto \zeta^g$ of $\mathbb{Q}(\zeta)$.
Then $\sigma$ generates $\operatorname{Gal}(K/\mathbb{Q})$, and $\sigma(\eta_i) = \eta_{i+1}$ for every $i \in \mathbb{Z}/N\mathbb{Z}$.

\begin{example} \label{ex:fields}
If $p^n = 3^4$, then $M = 243$ and $H = \{\pm 1\}$. We take $g = 2$, so that $\eta_i = 2\cos(2\pi \cdot 2^i/243)$.
If $p^n = 7^2$, then $M = 343$ and $H = \{\pm 1, \pm 18, \pm 19\}$. We take $g = 3$.
If $p^n = 29$, then $M = 841$ and $H$ has order $28$. We take $g = 2$.
\end{example}

\begin{lemma} \label{lem:field}
\begin{enumerate}
\item The discriminant of $K$ is $p^e$ with $e = \sum_{j=1}^{n} (j+1)(p-1)p^{j-1}$.
\item A prime number $q \ne p$ splits completely in $K$ if and only if $q^{p-1} \equiv 1 \pmod{M}$.
\item There is exactly one prime ideal of $K$ above $p$, and it is principal.
\item Let $c$ be an integer and let $(a_i)_{i \in \mathbb{Z}/N\mathbb{Z}}$ be a family of integers. Put $\alpha = c + \sum_i a_i \eta_i$. Then $\alpha \in O_K$ and
\[
N_{K/\mathbb{Q}}(\alpha) = \prod_{j \in \mathbb{Z}/N\mathbb{Z}} \Bigl( c + \sum_{i} a_i \eta_{i+j} \Bigr) .
\]
Moreover, $|N_{K/\mathbb{Q}}(\alpha)| \le b^N$ with $b = |c| + (p-1) \sum_i |a_i|$.
\end{enumerate}
\end{lemma}

\begin{proof}
(1) The group of Dirichlet characters associated with $K$ is cyclic of order $N$.
For $1 \le j \le n$, it contains exactly $(p-1)p^{j-1}$ characters of conductor $p^{j+1}$.
The assertion follows from the conductor-discriminant formula.

(2) The prime number $q$ splits completely in $K$ if and only if its Frobenius element in $\operatorname{Gal}(\mathbb{Q}(\zeta)/\mathbb{Q}) = (\mathbb{Z}/M\mathbb{Z})^\times$, which is the class of $q$, lies in $H$.

(3) The prime number $p$ is totally ramified in $\mathbb{Q}(\zeta)$, so that it is totally ramified in $K$.
Put $\pi = \prod_{a \in H} (1 - \zeta^a)$.
Then $\pi \in O_K$ and $N_{K/\mathbb{Q}}(\pi) = N_{\mathbb{Q}(\zeta)/\mathbb{Q}}(1 - \zeta) = p$.
Hence $\pi O_K$ is the prime ideal of $K$ above $p$.

(4) The first two assertions follow from the properties of $\sigma$ stated above.
Since $|\eta_i| \le p - 1$ for every $i$, each factor of the product has absolute value at most $b$.
\end{proof}

By \cref{lem:field}, the discriminants of $k_{3^4}$, $k_{7^2}$, and $k_{29}$ are $3^{364}$, $7^{138}$, and $29^{56}$, respectively.
For a number field $L$, let $\operatorname{rd}(L)$ denote the root discriminant of $L$.
Then
\begin{align*}
\log \operatorname{rd}(k_{3^4}) &= \tfrac{364}{81} \log 3 = 4.9369737416\ldots, \\
\log \operatorname{rd}(k_{7^2}) &= \tfrac{138}{49} \log 7 = 5.4803183789\ldots, \\
\log \operatorname{rd}(k_{29}) &= \tfrac{56}{29} \log 29 = 6.5023643613\ldots,
\end{align*}
so that $\operatorname{rd}(k_{3^4}) < 139.35$, $\operatorname{rd}(k_{7^2}) < 239.93$, and $\operatorname{rd}(k_{29}) < 666.73$.
A prime number $q \ne 3$ splits completely in $k_{3^4}$ if and only if $q \equiv \pm 1 \pmod{243}$.

\section{Three known results} \label{sec:known}

\subsection{An upper bound for class numbers}

Let $\gamma$ denote Euler's constant.
The following is \cite[Lemma 5.1]{Mi14}; see also \cite[Lemma 1]{Mi15b}.

\begin{proposition}[Miller] \label{prop:miller}
Let $L$ be a totally real number field of degree $N$, and let $c$ be a positive real number.
Define a function $F$ on $\mathbb{R}$ by
\[
F(x) = \frac{e^{-(x/c)^2}}{\cosh(x/2)} .
\]
Let $S$ be a set of prime numbers each of which splits completely in $L$ into principal prime ideals.
Put
\[
I(c) = \int_0^\infty \frac{1 - F(x)}{2} \left( \frac{1}{\sinh(x/2)} + \frac{1}{\cosh(x/2)} \right) dx
\]
and
\[
B = \frac{\pi}{2} + \gamma + \log 8\pi - \log \operatorname{rd}(L) - I(c) + 2 \sum_{q \in S} \sum_{m=1}^{\infty} \frac{\log q}{q^{m/2}} F(m \log q) .
\]
If $B > 0$, then the class number $h$ of $L$ satisfies
\[
h < \frac{2c\sqrt{\pi}}{N B} .
\]
\end{proposition}

The set $S$ may be any set of prime numbers with the stated property, and it may be empty.
Every term of the double sum is positive, so that a lower bound for $B$ is obtained if the sum is replaced by a partial sum.
Since $\cosh((m \log q)/2) = (q^{m/2} + q^{-m/2})/2$, the terms of the double sum satisfy
\begin{equation} \label{eq:term}
2 \cdot \frac{\log q}{q^{m/2}} F(m \log q) = \frac{4 \log q}{q^m + 1} \, e^{-(m \log q/c)^2} .
\end{equation}

\subsection{A divisibility result}

The following is a consequence of \cite[Theorems 6 and 8]{vdL82}; it is the general form of the rank corollary of Masley \cite[Corollary 2.15]{Ma78}.
We include a proof for the convenience of the reader.

\begin{proposition}[Masley, van der Linden] \label{prop:masley}
Let $E/F$ be a cyclic extension of number fields of degree $n$, and let $\ell$ be a prime number that does not divide $n$.
Suppose that $\ell$ does not divide the class number of any field $L$ with $F \subseteq L \subsetneq E$.
Let $f$ denote the order of $\ell$ in $(\mathbb{Z}/n\mathbb{Z})^\times$.
If $\ell$ divides the class number of $E$, then $\ell^f$ divides the class number of $E$.
\end{proposition}

\begin{proof}
Let $G$ denote the Galois group of $E/F$, let $\tau$ be a generator of $G$, and let $A$ denote the $\ell$-part of the ideal class group of $E$.
Since $\ell$ does not divide $n$, the ring $\mathbb{Z}_\ell[G]$ is the product of the rings $\mathbb{Z}_\ell[x]/(\Phi_d(x))$, where $d$ runs over the positive divisors of $n$ and $\Phi_d$ denotes the $d$-th cyclotomic polynomial.
Accordingly, $A$ is the direct sum of submodules $A_d$, where $\tau$ acts on $A_d$ through $\mathbb{Z}_\ell[x]/(\Phi_d(x))$.

Let $d$ be a positive divisor of $n$ with $d < n$, and let $L$ denote the subfield of $E$ fixed by $\tau^d$.
Then $F \subseteq L \subsetneq E$.
The element $\nu = \sum_{j=0}^{n/d-1} \tau^{dj}$ of $\mathbb{Z}_\ell[G]$ acts on $A$ as the composite of the norm map to the $\ell$-part of the class group of $L$ and the map induced by extension of ideals.
The $\ell$-part of the class group of $L$ is trivial by assumption, so that $\nu$ acts on $A$ as zero.
On the other hand, $\tau^d$ acts on $A_d$ as the identity, so that $\nu$ acts on $A_d$ as multiplication by $n/d$, which is a unit of $\mathbb{Z}_\ell$.
Hence $A_d = 0$.

Therefore $A = A_n$ is a module over $\mathbb{Z}_\ell[x]/(\Phi_n(x))$.
This ring is a product of copies of the ring of integers of the unramified extension of $\mathbb{Q}_\ell$ of degree $f$.
The order of every finite module over this ring is a power of $\ell^f$.
\end{proof}

\subsection{The class numbers of the subfields}

\begin{proposition}[Iwasawa] \label{prop:iwasawa}
Let $p$ be a prime number and let $n$ be a positive integer. Then $p$ does not divide $h(k_{p^n})$.
\end{proposition}

\begin{proof}
The extension $k_{p^n}/\mathbb{Q}$ is a Galois extension whose degree is a power of $p$, and $p$ is the only place of $\mathbb{Q}$ that ramifies in it.
Hence the assertion follows from a theorem of Iwasawa \cite{Iw56}; see \cite[Theorem 10.4]{Wa97}.
\end{proof}

\begin{proposition} \label{prop:sub}
The fields $k_3$, $k_{3^2}$, $k_{3^3}$, and $k_7$ have class number $1$.
\end{proposition}

\begin{proof}
These fields have the conductors $9$, $27$, $81$, and $49$.
By a theorem of van der Linden \cite[Theorem 1]{vdL82}, every real abelian field whose conductor is a prime power $q$ with $\varphi(q) \le 66$ has class number $1$.
This proves the assertion.

The assertion also follows from \cref{prop:miller,prop:masley,prop:iwasawa}, as we now explain.
We apply \cref{prop:miller} with the empty set $S$ and with the value of $c$ shown in \cref{tab:sub}.
The discriminants are given by \cref{lem:field}, and the bounds for $I(c)$ are obtained as in \cref{sec:bound}.
This yields the upper bounds for the class numbers in the last column of \cref{tab:sub}.

Since $h(k_3) < 2$ and $h(k_{3^2}) < 2$, these two class numbers are $1$.
Let $\ell$ be a prime number dividing $h(k_{3^3})$.
Then $\ell \ne 3$ by \cref{prop:iwasawa}.
By \cref{prop:masley}, applied to $E = k_{3^3}$ and $F = \mathbb{Q}$, the number $\ell^f$ divides $h(k_{3^3})$, where $f$ denotes the order of $\ell$ modulo $27$.
Here we use that $k_3$ and $k_{3^2}$ have class number $1$.
Since $h(k_{3^3}) < 6$, we have $\ell^f \le 5$, so that $\ell \in \{2, 5\}$ and $f \le 2$.
This is impossible, since $2$ and $5$ have order $18$ modulo $27$.
Hence $h(k_{3^3}) = 1$.
In the same way, a prime number $\ell$ dividing $h(k_7)$ satisfies $\ell \ne 7$ and $\ell^f \le 4$, where $f$ denotes the order of $\ell$ modulo $7$.
Then $\ell \in \{2, 3\}$ and $f \le 2$, which is impossible, since $2$ and $3$ have the orders $3$ and $6$ modulo $7$.
Hence $h(k_7) = 1$.
\end{proof}

\begin{table}[h]
\centering
\begin{tabular}{lrcrccc}
\toprule
field & degree & $\log \operatorname{rd} \le$ & $c$ & $I(c) \le$ & $B \ge$ & $h <$ \\
\midrule
$k_3$ & $3$ & $1.4648164$ & $2$ & $2.1508333$ & $1.7565338$ & $1.3455$ \\
$k_{3^2}$ & $9$ & $2.6854968$ & $3$ & $1.7991266$ & $0.8875601$ & $1.3314$ \\
$k_{3^3}$ & $27$ & $3.8247984$ & $8$ & $1.3711535$ & $0.1762316$ & $5.9601$ \\
$k_7$ & $7$ & $3.3358460$ & $5$ & $1.5090534$ & $0.5272841$ & $4.8022$ \\
\bottomrule
\end{tabular}
\caption{Bounds obtained from \cref{prop:miller} with the empty set $S$. Every entry is rounded in the safe direction, and the last two columns are computed from the values before rounding.}
\label{tab:sub}
\end{table}

\begin{corollary} \label{cor:threshold}
\begin{enumerate}
\item If $h(k_{3^4}) > 1$, then $h(k_{3^4}) \ge 163$.
\item If $h(k_{7^2}) > 1$, then $h(k_{7^2}) \ge 197$.
\item If $h(k_{29}) > 1$, then $h(k_{29}) \ge 59$.
\end{enumerate}
\end{corollary}

\begin{proof}
Let $p^n$ be $3^4$, $7^2$, or $29$, and let $\ell$ be a prime number dividing $h(k_{p^n})$.
Then $\ell \ne p$ by \cref{prop:iwasawa}.
Every field $L$ with $\mathbb{Q} \subseteq L \subsetneq k_{p^n}$ has class number $1$: this holds by \cref{prop:sub} if $p^n$ is $3^4$ or $7^2$, and the only such field is $\mathbb{Q}$ if $p^n = 29$.
By \cref{prop:masley}, applied to $E = k_{p^n}$ and $F = \mathbb{Q}$, the number $\ell^f$ divides $h(k_{p^n})$, where $f$ denotes the order of $\ell$ modulo $p^n$.

(1) We show that $\ell^f \ge 163$.
If $f = 1$, then $\ell \equiv 1 \pmod{81}$, and the least prime number with this property is $163$.
If $f \ge 2$ and $\ell^f < 163$, then $\ell \in \{2, 5, 7, 11\}$.
The orders of $2$, $5$, $7$, and $11$ modulo $81$ are $54$, $54$, $27$, and $54$, respectively, so that $\ell^f > 163$ in each of these cases.

(2) We show that $\ell^f \ge 197$.
If $f = 1$, then $\ell \equiv 1 \pmod{49}$, and the least prime number with this property is $197$.
If $f \ge 2$ and $\ell^f < 197$, then $\ell \in \{2, 3, 5, 11, 13\}$.
The orders of $2$, $3$, $5$, $11$, and $13$ modulo $49$ are $21$, $42$, $42$, $21$, and $14$, respectively, so that $\ell^f > 197$ in each of these cases.

(3) We show that $\ell^f \ge 59$.
If $f = 1$, then $\ell \equiv 1 \pmod{29}$, and the least prime number with this property is $59$.
If $f \ge 2$ and $\ell^f < 59$, then $\ell \in \{2, 3, 5, 7\}$.
The orders of $2$, $3$, $5$, and $7$ modulo $29$ are $28$, $28$, $14$, and $7$, respectively, so that $\ell^f > 59$ in each of these cases.
\end{proof}

\section{Principal prime ideals from norms} \label{sec:relations}

Throughout this section, $K = k_{p^n}$ is as in \cref{sec:field}, and we assume that every field $k_{p^j}$ with $0 \le j < n$ has class number $1$.
By \cref{prop:sub}, this assumption holds for $p^n = 3^4$ and for $p^n = 7^2$; it holds for $p^n = 29$ since $k_1 = \mathbb{Q}$.
We say that a prime number $r$ is \emph{settled} if every prime ideal of $K$ above $r$ is principal.
The prime number $p$ is settled by \cref{lem:field}.

\begin{lemma} \label{lem:nonsplit}
Let $r \ne p$ be a prime number that does not split completely in $K$. Then $r$ is settled.
\end{lemma}

\begin{proof}
Let $D$ denote the decomposition field of $r$ in $K$.
Then $D = k_{p^j}$ for some $j < n$, so that $h(D) = 1$.
Let $\mathfrak{r}$ be a prime ideal of $D$ above $r$.
The extension $K/D$ is cyclic, the prime ideal $\mathfrak{r}$ is unramified in it, and the decomposition group of a prime ideal of $K$ above $\mathfrak{r}$ is $\operatorname{Gal}(K/D)$.
Hence $\mathfrak{r} O_K$ is a prime ideal of $K$.
It is principal since $\mathfrak{r}$ is principal.
Every prime ideal of $K$ above $r$ is of this form.
\end{proof}

\begin{definition} \label{def:relation}
Let $q$ be a prime number that splits completely in $K$.
A \emph{relation for $q$} is an element $\alpha \in O_K$ such that $q$ divides the integer $N_{K/\mathbb{Q}}(\alpha)$ and $q^2$ does not divide it.
\end{definition}

\begin{lemma} \label{lem:relation}
Let $q$ be a prime number that splits completely in $K$, and let $\alpha$ be a relation for $q$.
Suppose that every prime number $r \ne q$ dividing $N_{K/\mathbb{Q}}(\alpha)$ is settled.
Then $q$ is settled.
\end{lemma}

\begin{proof}
Every prime ideal of $K$ above $q$ has norm $q$.
Since $q$ divides $N_{K/\mathbb{Q}}(\alpha)$ exactly once, there is exactly one prime ideal $\mathfrak{q}$ of $K$ above $q$ that contains $\alpha$, and $\alpha O_K = \mathfrak{q} \mathfrak{a}$, where $\mathfrak{a}$ is a product of prime ideals above prime numbers $r \ne q$ dividing $N_{K/\mathbb{Q}}(\alpha)$.
The ideal $\mathfrak{a}$ is principal by assumption.
Hence $\mathfrak{q}$ is principal.
The extension $K/\mathbb{Q}$ is Galois, so that every prime ideal of $K$ above $q$ is of the form $\tau(\mathfrak{q})$ with $\tau \in \operatorname{Gal}(K/\mathbb{Q})$.
\end{proof}

By \cref{lem:nonsplit}, the assumption of \cref{lem:relation} only concerns the prime numbers $r \ne q$ that divide the norm and split completely in $K$.
We use \cref{lem:relation} in two ways.

\begin{proposition} \label{prop:list}
Let $q_1, \ldots, q_t$ be prime numbers that split completely in $K$, and let $T$ be a set of settled prime numbers.
For $1 \le i \le t$, let $\alpha_i$ be a relation for $q_i$.
Suppose that, for every $i$, every prime number $r \ne q_i$ that divides $N_{K/\mathbb{Q}}(\alpha_i)$ and splits completely in $K$ belongs to $T \cup \{q_1, \ldots, q_{i-1}\}$.
Then $q_1, \ldots, q_t$ are settled.
\end{proposition}

\begin{proof}
This follows from \cref{lem:nonsplit,lem:relation} by induction on $i$.
\end{proof}

\begin{proposition} \label{prop:range}
Let $A < X$ be positive real numbers.
Suppose that every prime number $q < A$ that splits completely in $K$ is settled.
Suppose further that, for every prime number $q$ with $A \le q < X$ that splits completely in $K$, there is a relation $\alpha_q$ for $q$ such that every prime number $r \ne q$ that divides $N_{K/\mathbb{Q}}(\alpha_q)$ and splits completely in $K$ satisfies $r < q$.
Then every prime number $q < X$ that splits completely in $K$ is settled.
\end{proposition}

\begin{proof}
Suppose that the assertion is false, and let $q$ be the least prime number $q < X$ that splits completely in $K$ and is not settled.
Then $q \ge A$.
Every prime number $r \ne q$ that divides $N_{K/\mathbb{Q}}(\alpha_q)$ either does not split completely in $K$, or equals $p$, or splits completely in $K$ and satisfies $r < q$.
In each case $r$ is settled, by \cref{lem:nonsplit}, by \cref{lem:field}, and by the choice of $q$, respectively.
Hence $q$ is settled by \cref{lem:relation}, a contradiction.
\end{proof}

In \cref{prop:range}, the condition on $q$ does not refer to the other relations.
Hence it can be verified separately for disjoint intervals of prime numbers $q$.

In $k_{7^2}$, the least prime number that splits completely is $19$, and it divides the norms of most elements of small height.
For this prime number we use the following lemma.

\begin{lemma} \label{lem:block}
Let $q_0$ be a prime number that splits completely in $K$, and let $\mathfrak{q}_1, \ldots, \mathfrak{q}_N$ denote the prime ideals of $K$ above $q_0$.
Let $\beta_1, \ldots, \beta_s$ be elements of $O_K$ with the following two properties for every $i$.
\begin{enumerate}
\item We have $|N_{K/\mathbb{Q}}(\beta_i)| = p^{a_i} q_0^{k_i}$ with integers $a_i \ge 0$ and $k_i \ge 1$.
\item Exactly $k_i$ of the prime ideals $\mathfrak{q}_1, \ldots, \mathfrak{q}_N$ contain $\beta_i$.
\end{enumerate}
For $\beta \in O_K$, let $v(\beta) \in \{0,1\}^N$ denote the vector whose $j$-th entry is $1$ if and only if $\beta \in \mathfrak{q}_j$.
If the vectors $v(\tau(\beta_i))$ with $1 \le i \le s$ and $\tau \in \operatorname{Gal}(K/\mathbb{Q})$ generate the group $\mathbb{Z}^N$, then $q_0$ is settled.
\end{lemma}

\begin{proof}
Let $\beta$ be one of the elements $\tau(\beta_i)$.
By (1) and (2), the valuation of $\beta$ at $\mathfrak{q}_j$ is the $j$-th entry of $v(\beta)$, and $\beta O_K$ is the product of the ideals $\mathfrak{q}_j$ with $\beta \in \mathfrak{q}_j$ and a power of the prime ideal above $p$.
The latter is principal by \cref{lem:field}.
Hence $v(\beta)$ lies in the kernel of the homomorphism from $\mathbb{Z}^N$ to the ideal class group of $K$ that sends the $j$-th standard basis vector to the class of $\mathfrak{q}_j$.
If these vectors generate $\mathbb{Z}^N$, then this homomorphism is trivial.
\end{proof}

\begin{remark} \label{rem:blockcheck}
The conditions of \cref{lem:block} are verified as follows.
Suppose that $\Phi_M$ is the product of $N$ distinct irreducible polynomials $g_1, \ldots, g_N$ modulo $q_0$.
Then the ideals $(q_0, g_j(\zeta))$ are the prime ideals of $\mathbb{Z}[\zeta]$ above $q_0$.
There are $N$ of them, so that each prime ideal $\mathfrak{q}_i$ lies below exactly one of them.
Let $\beta \in O_K$ be represented by a polynomial $b \in \mathbb{Z}[x]$ with $b(\zeta) = \beta$.
After renumbering, we have $\beta \in \mathfrak{q}_j$ if and only if $g_j$ divides $b$ modulo $q_0$.
\end{remark}

\section{The lists of relations} \label{sec:data}

In this section, a prime number is called a \emph{split prime} of $K$ if it splits completely in $K$.
All elements are given in the form $c + \sum_i a_i \eta_i$ of \cref{lem:field}, with the generators $g$ of \cref{ex:fields}.

\subsection{The field $k_{3^4}$}

Let $K = k_{3^4}$.
Let $S_3$ denote the set of prime numbers $q < 3 \cdot 10^{10}$ with $q \equiv \pm 1 \pmod{243}$.
The set $S_3$ has $16049777$ elements.
Its least elements are $487$, $971$, $1459$, and $2917$, and its greatest element is $29999999969$.

\begin{proposition} \label{prop:data3}
Every prime number in $S_3$ is settled; that is, every prime ideal of $k_{3^4}$ above a prime number in $S_3$ is principal.
\end{proposition}

\begin{proof}
The accompanying data consist of two parts.

The first part is a list of $8699$ relations in the order required by \cref{prop:list}, with the empty set $T$.
The prime numbers $q_i$ of this list include all $8160$ elements of $S_3$ below $10^7$.
The remaining $539$ prime numbers $q_i$ are split primes greater than $10^7$ that occur in the norms of later relations.
By \cref{prop:list}, every element of $S_3$ below $10^7$ is settled.

The second part gives, for each of the $16041617$ elements $q$ of $S_3$ with $q \ge 10^7$, a relation $\alpha_q$ for $q$ such that every split prime $r \ne q$ dividing $N_{K/\mathbb{Q}}(\alpha_q)$ satisfies $r < q$.
By \cref{prop:range} with $A = 10^7$ and $X = 3 \cdot 10^{10}$, every element of $S_3$ is settled.
The verification of the data is described in \cref{subsec:verification}.
\end{proof}

\begin{example} \label{ex:relations}
The relation for the greatest element $q = 29999999969$ of $S_3$ is
\[
\alpha = \eta_1 + \eta_3 - \eta_8 + \eta_{10} - \eta_{17} + \eta_{20} - \eta_{22} - \eta_{23} - \eta_{29} + \eta_{37},
\]
and $|N_{K/\mathbb{Q}}(\alpha)| = 3^2 \cdot 487 \cdot 1459 \cdot 2917 \cdot 23849479 \cdot q$.
The relation for $q = 8742169$ in the first part is $\alpha = -1 - \eta_{11} - \eta_{14} - \eta_{23} + \eta_{29}$, and $|N_{K/\mathbb{Q}}(\alpha)| = q$.
\end{example}

\subsection{The field $k_{7^2}$}

Let $K = k_{7^2}$.
The least split primes of $K$ are $19$, $1373$, $1697$, and $1733$, and there are $28381$ split primes $q \le 2.2 \cdot 10^7$.

\begin{proposition} \label{prop:data7}
There is a set $S_7$ of $29835$ split primes of $k_{7^2}$ such that $S_7$ contains every split prime $q \le 2.2 \cdot 10^7$ and every element of $S_7$ is settled.
\end{proposition}

\begin{proof}
The prime number $19$ is settled by \cref{lem:block}, applied to $q_0 = 19$ and to the eight elements
\begin{gather*}
\pm(2 + \eta_0 + \eta_{24}), \quad \pm(\eta_0 - \eta_{26}), \quad \eta_0 - \eta_{21}, \\
\eta_{23} - \eta_0, \quad \eta_{28} - \eta_0, \quad \eta_{46} - \eta_0 - 1,
\end{gather*}
whose norms have the absolute values $7 \cdot 19^4$, $7^2 \cdot 19^4$, $7^9 \cdot 19^2$, $7^2 \cdot 19^4$, $7^9 \cdot 19^2$, and $19^3$, respectively.
The order of $19$ modulo $343$ is $6$, so that $\Phi_{343}$ is the product of $49$ irreducible polynomials of degree $6$ modulo $19$, and the conditions are verified as in \cref{rem:blockcheck}.
The $8 \cdot 49$ vectors generate $\mathbb{Z}^{49}$; this is verified by computing the Hermite normal form of the matrix formed by these vectors.

The accompanying data then give a list of $29834$ relations in the order required by \cref{prop:list}, with $T = \{19\}$.
The prime numbers $q_i$ of this list, together with $19$, form the set $S_7$.
The greatest element of $S_7$ is $59529496693$.
\end{proof}

\subsection{The field $k_{29}$}

Let $K = k_{29}$.
Let $S_{29}$ denote the set of prime numbers $q < 4 \cdot 10^9$ with $q^{28} \equiv 1 \pmod{841}$.
The set $S_{29}$ has $6549261$ elements.
Its least elements are $41$, $137$, $467$, and $571$, and its greatest element is $3999999659$.

\begin{proposition} \label{prop:data29}
Every prime number in $S_{29}$ is settled.
\end{proposition}

\begin{proof}
The accompanying data consist of two parts.
The first part gives, for each of the $327$ elements $q$ of $S_{29}$ below $10^5$, an element $\alpha_q \in O_K$ with $|N_{K/\mathbb{Q}}(\alpha_q)| = q$.
Such an element is a relation for $q$ whose norm has no other prime factor, so that $q$ is settled by \cref{lem:relation}.
The second part gives, for each of the $6548934$ elements $q$ of $S_{29}$ with $q \ge 10^5$, a relation $\alpha_q$ for $q$ such that every split prime $r \ne q$ dividing $N_{K/\mathbb{Q}}(\alpha_q)$ satisfies $r < q$.
By \cref{prop:range} with $A = 10^5$ and $X = 4 \cdot 10^9$, every element of $S_{29}$ is settled.
\end{proof}

\begin{remark} \label{rem:41}
In $k_{29}$, we did not find elements to which \cref{lem:block} applies with $q_0 = 41$.
The elements $\alpha_q$ of the first part were computed with the functions \texttt{bnfinit} and \texttt{bnfisprincipal} of PARI/GP \cite{PARI}, and then written in the form $c + \sum_i a_i \eta_i$.
The integers $c$ and $a_i$ are larger than in the other relations: they have at most four digits for $234$ of the $327$ prime numbers, and at most $24$ digits in all cases.
The function \texttt{bnfinit} assumes the generalized Riemann hypothesis.
This does not affect the proof, since the proof only uses the elements $\alpha_q$ themselves, and their norms are verified as described below.
For example, the relation for $q = 499979$ in the second part is $\alpha = 1 + \eta_1 + \eta_3 + \eta_4 + \eta_7 - \eta_{11} + \eta_{12}$, and $|N_{K/\mathbb{Q}}(\alpha)| = 41^3 \cdot 202621 \cdot 271879 \cdot q$.
\end{remark}

\subsection{The verification} \label{subsec:verification}

A relation is stored as the prime number $q$, the integers $c$ and $a_i$, and the factorization of $|N_{K/\mathbb{Q}}(\alpha)|$ into prime numbers.
The verification program reads these data and checks the following assertions with exact integer arithmetic.
\begin{enumerate}
\item Every number in the stated factorization is a prime number.
\item The product of the stated factorization equals $|N_{K/\mathbb{Q}}(\alpha)|$.
\item The prime number $q$ splits completely in $K$, and its exponent in the factorization is $1$.
\item The split primes $r \ne q$ in the factorization satisfy the condition of \cref{prop:list} or of \cref{prop:range}.
\item Every split prime of $K$ in the interval under consideration occurs as the prime number $q$ of some relation.
\end{enumerate}
For (1) and (5), primality is proved by the function \texttt{fmpz\_is\_prime} of FLINT \cite{FLINT}; all prime numbers that occur are less than $2^{36}$.
For (5), the program tests every integer in the interval that satisfies the congruence of \cref{lem:field}.

For (2), the norm is computed as follows.
Let $l$ be a prime number with $l \equiv 1 \pmod{M}$, and let $z \in \mathbb{F}_l$ be an element of order $M$.
There is a ring homomorphism from $\mathbb{Z}[\zeta]$ to $\mathbb{F}_l$ that sends $\zeta$ to $z$.
By \cref{lem:field}, the image of $N_{K/\mathbb{Q}}(\alpha)$ under this homomorphism is
\[
\prod_{j \in \mathbb{Z}/N\mathbb{Z}} \Bigl( c + \sum_i a_i \rho_{i+j} \Bigr), \quad \text{where} \quad \rho_i = \sum_{a \in H} z^{g^i a} .
\]
The program uses twelve such prime numbers $l > 2^{61}$, and it adds further ones if they are needed for the condition below.
By \cref{lem:field}, the integer $N_{K/\mathbb{Q}}(\alpha)$ has absolute value at most $b^N$, where $b = |c| + (p-1)\sum_i |a_i|$.
It is therefore determined by its residues modulo prime numbers $l$ whose product exceeds $2b^N + 1$, and the program checks that this condition holds.

\begin{remark} \label{rem:search}
The relations were found as follows; this is not part of the proof.
Let $q$ be a split prime of $K$.
The elements $c + \sum_i a_i \eta_i$ that lie in a fixed prime ideal above $q$ correspond to the vectors $(c, (a_i)_i)$ of a lattice of index $q$, which is described by the images of the Gaussian periods in $\mathbb{F}_q$.
We reduce a basis of this lattice with the LLL algorithm, compute the norms of the short vectors and of their sums and differences, and factor these norms with PARI/GP \cite{PARI} until a relation with the required property is found.
This takes about $20$ to $35$ milliseconds per prime number on one core.
For $k_{7^2}$, the whole search took about $20$ minutes on two cores, and the verification takes less than $15$ seconds.
For $k_{3^4}$, the interval from $10^7$ to $3 \cdot 10^{10}$ was divided into $82$ intervals, and the search and the verification took between $16$ and $41$ minutes per interval on machines with four cores.
For $k_{29}$, the interval from $5 \cdot 10^5$ to $4 \cdot 10^9$ was divided into $80$ intervals, and the search and the verification took between $2$ and $7$ minutes per interval.
The norms of the elements of small height depend on the field.
In our samples of elements with $|c| \le 2$ and at most five nonzero coefficients $a_i \in \{\pm 1\}$, typical norms were of the order of $10^{12}$ in $k_{3^4}$, $10^{18}$ in $k_{7^2}$, and $10^{20}$ in $k_{29}$; these figures depend on the sample and are only an indication.
\end{remark}

\begin{remark} \label{rem:second}
All relations were verified a second time with an independent program.
It proves primality with the function \texttt{isprime} of PARI/GP.
For $k_{3^4}$ and $k_{7^2}$, it computes the resultant of $\Phi_M$ and the polynomial $c + \sum_i a_i \sum_{a \in H} x^{g^i a \bmod M}$ with PARI/GP; its absolute value is $|N_{K/\mathbb{Q}}(\alpha)|^{p-1}$.
For $k_{29}$, this resultant is slow, and the program computes the norm as the determinant of the matrix of the multiplication by $\alpha$ with respect to the basis $1, \eta_0, \ldots, \eta_{27}$ of $K$ over $\mathbb{Q}$; the resultant is used for about one relation in $800$.
\end{remark}

\section{The upper bounds} \label{sec:bound}

\begin{lemma} \label{lem:integral}
Let $c$, $\varepsilon$, and $R$ be positive real numbers with $\varepsilon < R$, and let $u$ denote the integrand in the definition of $I(c)$.
Then
\[
I(c) \le \int_\varepsilon^R u(x)\, dx + \left( \frac{1}{c^2} + \frac{1}{8} \right) \left( \frac{\varepsilon^2}{2} + \frac{\varepsilon^3}{6} \right) + 2 e^{-R/2} \left( 1 + \frac{1}{1 - e^{-R}} \right).
\]
\end{lemma}

\begin{proof}
Let $x$ be a positive real number.
From $\cosh(y) \le e^{y^2/2}$ and $e^{-t} \ge 1 - t$, we obtain
\[
F(x) \ge e^{-x^2/c^2 - x^2/8} \ge 1 - \frac{x^2}{c^2} - \frac{x^2}{8} .
\]
Together with $1/\sinh(x/2) \le 2/x$ and $1/\cosh(x/2) \le 1$, this gives
\[
u(x) \le \left( \frac{1}{c^2} + \frac{1}{8} \right) \left( x + \frac{x^2}{2} \right),
\]
and the integral of the right-hand side over the interval from $0$ to $\varepsilon$ is the second term of the assertion.

Suppose that $x \ge R$.
Then $0 \le 1 - F(x) \le 1$, $1/\sinh(x/2) = 2e^{-x/2}/(1 - e^{-x}) \le 2e^{-x/2}/(1 - e^{-R})$, and $1/\cosh(x/2) \le 2e^{-x/2}$.
Hence
\[
u(x) \le e^{-x/2} \left( 1 + \frac{1}{1 - e^{-R}} \right),
\]
and the integral of the right-hand side over the interval from $R$ to $\infty$ is the third term of the assertion.
\end{proof}

For a set $S$ of prime numbers and a positive real number $c$, let $\Sigma(S, c)$ denote the double sum in \cref{prop:miller}.

\begin{proposition} \label{prop:bound}
\begin{enumerate}
\item Let $c = 100$. Then $I(c) \le 1.2646927462$ and $\Sigma(S_3, c) \ge 0.8733014006$. For $L = k_{3^4}$ and $S = S_3$, the number $B$ of \cref{prop:miller} satisfies $B \ge 0.0438183$, and $h(k_{3^4}) < 99.8767$.
\item Let $c = 80$. Then $I(c) \le 1.2651137825$ and $\Sigma(S_7, c) \ge 1.4113091085$. For $L = k_{7^2}$ and $S = S_7$, the number $B$ of \cref{prop:miller} satisfies $B \ge 0.0380603$, and $h(k_{7^2}) < 152.0639$.
\item Let $c = 80$. Then $\Sigma(S_{29}, c) \ge 2.6238431599$. For $L = k_{29}$ and $S = S_{29}$, the number $B$ of \cref{prop:miller} satisfies $B \ge 0.2285484$, and $h(k_{29}) < 42.7877$.
\end{enumerate}
\end{proposition}

\begin{proof}
The hypothesis of \cref{prop:miller} on the sets $S_3$, $S_7$, and $S_{29}$ holds by \cref{prop:data3,prop:data7,prop:data29}.
The inequalities are obtained with ball arithmetic, using the library Arb \cite{Jo17}, which is now part of FLINT \cite{FLINT}.
For the integral, we use \cref{lem:integral} with $\varepsilon = 10^{-3}$ and $R = 200$, and the integral over the interval from $\varepsilon$ to $R$ is enclosed by the rigorous integration algorithm of \cite{Jo18}.
For the sum, we use a partial sum, which is a lower bound since every term is positive.
In (2), the partial sum consists of the terms with $m \le 8$.
In (1) and (3), the terms are computed by \eqref{eq:term}, and the partial sum consists of the terms with $m \le 8$ and $q^{m-1} \le 10^{16}$.
The sets $S_3$ and $S_{29}$ are computed again in this step, and their elements are proved to be prime numbers.

We have $\pi/2 + \gamma + \log 8\pi = 5.3721834192\ldots$.
With the values of $\log \operatorname{rd}$ of \cref{sec:field}, we obtain in (1)
\begin{align*}
B &\ge 5.3721834192 - 4.9369737417 - 1.2646927462 + 0.8733014006 \\
&> 0.0438183,
\end{align*}
and in (2)
\begin{align*}
B &\ge 5.3721834192 - 5.4803183790 - 1.2651137825 + 1.4113091085 \\
&> 0.0380603,
\end{align*}
and in (3)
\begin{align*}
B &\ge 5.3721834192 - 6.5023643614 - 1.2651137825 + 2.6238431599 \\
&> 0.2285484 .
\end{align*}
The bounds for the class numbers follow from \cref{prop:miller} with $N = 81$, $N = 49$, and $N = 29$; they are computed from the values of $B$ before rounding.
\end{proof}

\begin{table}[h]
\centering
\begin{tabular}{lrcccc}
\toprule
field & $c$ & $I(c) \le$ & $\Sigma \ge$ & $B \ge$ & $h <$ \\
\midrule
$k_{3^4}$ & $60$ & $1.2660224223$ & $0.8355826772$ & $0.0047699$ & $550.5021$ \\
 & $80$ & $1.2651137825$ & $0.8611113518$ & $0.0312072$ & $112.1901$ \\
 & $100$ & $1.2646927462$ & $0.8733014006$ & $0.0438183$ & $99.8767$ \\
 & $150$ & $1.2642766179$ & $0.8855807997$ & $0.0565138$ & $116.1599$ \\
\midrule
$k_{7^2}$ & $60$ & $1.2660224223$ & $1.3962035757$ & $0.0220461$ & $196.8913$ \\
 & $80$ & $1.2651137825$ & $1.4113091085$ & $0.0380603$ & $152.0639$ \\
 & $100$ & $1.2646927462$ & $1.4184096767$ & $0.0455819$ & $158.7142$ \\
 & $150$ & $1.2642766179$ & $1.4254912948$ & $0.0530797$ & $204.4427$ \\
\midrule
$k_{29}$ & $60$ & $1.2660224223$ & $2.5675506639$ & $0.1713472$ & $42.8037$ \\
 & $80$ & $1.2651137825$ & $2.6238431599$ & $0.2285484$ & $42.7877$ \\
 & $100$ & $1.2646927462$ & $2.6505824088$ & $0.2557087$ & $47.8037$ \\
 & $150$ & $1.2642766179$ & $2.6774279523$ & $0.2829703$ & $64.7974$ \\
\bottomrule
\end{tabular}
\caption{Bounds obtained from \cref{prop:miller} with the sets $S_3$, $S_7$, and $S_{29}$. The column $\Sigma$ is a partial sum of the double sum. Every entry is rounded in the safe direction, and the last two columns are computed from the values before rounding.}
\label{tab:bounds}
\end{table}

\begin{proof}[Proof of \cref{thm:main}]
By \cref{prop:bound}, we have $h(k_{3^4}) < 99.8767$, $h(k_{7^2}) < 152.0639$, and $h(k_{29}) < 42.7877$.
By \cref{cor:threshold}, each of these class numbers is either $1$ or at least $163$, $197$, and $59$, respectively.
Hence the three class numbers are $1$.
\end{proof}

\section{Three first layers under the generalized Riemann hypothesis} \label{sec:grh}

Throughout this section, $n = 1$, so that $K = k_p$, $N = p$, and $M = p^2$.
The only proper subfield of $K$ is $\mathbb{Q}$, so that the assumption of \cref{sec:relations} holds.
By a \emph{split prime} we mean a prime number that splits completely in $K$.
By \cref{lem:nonsplit}, every prime number $r \ne p$ that is not a split prime is settled; such a prime number generates a prime ideal of $K$.

In \cref{sec:relations}, the relations are used one after another, and the first relations must be elements whose norms contain no split prime that is not yet settled.
For the fields of this section we could not compute such elements, because the norms of the elements of small height are large; see \cref{rem:searchgrh}.
Instead, we use many elements at the same time: the factorizations of the ideals that they generate are regarded as vectors over a cyclotomic ring, and we show by linear algebra that these vectors generate a free module.

Put $s = \sigma^{-1}$ and
\[
A = \mathbb{Z}[x]/(1 + x + \dots + x^{p-1}) .
\]
Let $\omega$ denote the complex number $e^{2\pi i/p}$.
The ring homomorphism $A \to \mathbb{Z}[\omega]$ that sends $x$ to $\omega$ is an isomorphism.
For $a \in A$, let $\operatorname{Nm}(a)$ denote the norm of the image of $a$ with respect to $\mathbb{Q}(\omega)/\mathbb{Q}$; it is an integer.
For a finite set $T$, let $A^T$ denote the free $A$-module with basis $(e_q)_{q \in T}$.

\subsection{The class group as a module}

Let $\operatorname{Cl}$ denote the ideal class group of $K$, written additively, and let $[\mathfrak{a}]$ denote the class of a nonzero ideal $\mathfrak{a}$ of $O_K$.

\begin{lemma} \label{lem:module}
There is a unique structure of $A$-module on $\operatorname{Cl}$ such that $x \cdot [\mathfrak{a}] = [s(\mathfrak{a})]$ for every nonzero ideal $\mathfrak{a}$ of $O_K$.
\end{lemma}

\begin{proof}
The group $\operatorname{Gal}(K/\mathbb{Q})$ is cyclic of order $p$ and generated by $s$, so that $\operatorname{Cl}$ is a module over $\mathbb{Z}[x]/(x^p - 1)$, where $x$ acts as $s$.
For every nonzero ideal $\mathfrak{a}$ of $O_K$, the ideal $\prod_{j=0}^{p-1} s^j(\mathfrak{a})$ is generated by the positive integer $N_{K/\mathbb{Q}}(\mathfrak{a})$.
Hence $1 + x + \dots + x^{p-1}$ acts on $\operatorname{Cl}$ as zero.
\end{proof}

For each split prime $q$, we fix a prime ideal $\mathfrak{q}_q$ of $K$ above $q$.
The prime ideals of $K$ above $q$ are the ideals $s^j(\mathfrak{q}_q)$ with $j \in \mathbb{Z}/p\mathbb{Z}$, and they are pairwise distinct.
Let $\mathfrak{p}$ denote the prime ideal of $K$ above $p$.

\begin{lemma} \label{lem:vector}
Let $\alpha$ be a nonzero element of $O_K$.
Suppose that
\[
|N_{K/\mathbb{Q}}(\alpha)| = p^a \prod_{q \in Q} q^{k_q},
\]
where $a$ is a nonnegative integer, $Q$ is a finite set of split primes, and the $k_q$ are positive integers.
For $q \in Q$, put
\[
J_q = \{ j \in \mathbb{Z}/p\mathbb{Z} \mid \alpha \in s^j(\mathfrak{q}_q) \} .
\]
If the set $J_q$ has exactly $k_q$ elements for every $q \in Q$, then
\[
\alpha O_K = \mathfrak{p}^a \prod_{q \in Q} \prod_{j \in J_q} s^j(\mathfrak{q}_q)
\qquad \text{and} \qquad
\sum_{q \in Q} \Bigl( \sum_{j \in J_q} x^j \Bigr) \cdot [\mathfrak{q}_q] = 0 \ \text{ in } \operatorname{Cl} .
\]
\end{lemma}

\begin{proof}
Let $q \in Q$, and for $j \in \mathbb{Z}/p\mathbb{Z}$ let $v_j$ denote the exponent of $s^j(\mathfrak{q}_q)$ in the factorization of $\alpha O_K$.
Every prime ideal above $q$ has norm $q$, so that $\sum_j v_j = k_q$.
We have $v_j \ge 1$ if $j \in J_q$ and $v_j = 0$ otherwise.
Since $J_q$ has $k_q$ elements, it follows that $v_j = 1$ for every $j \in J_q$.
The only prime ideals that contain $\alpha$ are $\mathfrak{p}$, if $a \ge 1$, and prime ideals above the prime numbers in $Q$.
This proves the first equality.
The second equality follows from \cref{lem:module}, because $\alpha O_K$ and $\mathfrak{p}$ are principal ideals; see \cref{lem:field}.
\end{proof}

\begin{definition} \label{def:vector}
Let $T$ be a finite set of split primes.
An element $\alpha \in O_K$ is \emph{admissible for $T$} if it satisfies the assumptions of \cref{lem:vector} with a set $Q \subseteq T$.
In this case, the \emph{vector of $\alpha$} is the element
\[
v(\alpha) = \sum_{q \in Q} \Bigl( \sum_{j \in J_q} x^j \Bigr) e_q
\]
of $A^T$.
\end{definition}

\begin{proposition} \label{prop:span}
Let $T$ be a finite set of split primes, and let $\alpha_1, \dots, \alpha_r$ be elements of $O_K$ that are admissible for $T$.
If the vectors $v(\alpha_1), \dots, v(\alpha_r)$ generate the $A$-module $A^T$, then every prime number in $T$ is settled.
\end{proposition}

\begin{proof}
Let $\varphi \colon A^T \to \operatorname{Cl}$ denote the homomorphism of $A$-modules that sends $e_q$ to $[\mathfrak{q}_q]$ for every $q \in T$.
By \cref{lem:vector}, we have $\varphi(v(\alpha_i)) = 0$ for every $i$.
If the vectors generate $A^T$, then $\varphi = 0$, so that $\mathfrak{q}_q$ is principal for every $q \in T$.
Then the conjugates of $\mathfrak{q}_q$ are principal as well.
\end{proof}

\subsection{Two criteria for generation}

\begin{lemma} \label{lem:pivot}
Let $T$ be a finite set, let $V$ be a finite subset of $A^T$, and let $q_0 \in T$ and $v_0 \in V$.
Suppose that the coefficient $u$ of $e_{q_0}$ in $v_0$ is a unit of $A$.
For each $v \in V$ with $v \ne v_0$, let $b_v$ denote the coefficient of $e_{q_0}$ in $v$, and put $v' = v - b_v u^{-1} v_0$; the coefficient of $e_{q_0}$ in $v'$ is zero, and we regard $v'$ as an element of $A^{T \setminus \{q_0\}}$.
If the elements $v'$ generate $A^{T \setminus \{q_0\}}$, then $V$ generates $A^T$.
\end{lemma}

\begin{proof}
Let $\Lambda$ denote the submodule of $A^T$ generated by $V$.
Every $v'$ belongs to $\Lambda$.
If the elements $v'$ generate $A^{T \setminus \{q_0\}}$, then $e_q \in \Lambda$ for every $q \in T$ with $q \ne q_0$.
Then $u e_{q_0} \in \Lambda$, because $v_0 - u e_{q_0}$ is a linear combination of the elements $e_q$ with $q \ne q_0$.
Since $u$ is a unit, we obtain $e_{q_0} \in \Lambda$.
\end{proof}

The elements $\pm x^j$ of $A$ are units, and \cref{lem:pivot} will be applied only when $u$ is of this form.
If a subset of $V$ generates $A^T$, then $V$ generates $A^T$; hence elements of $V$ may be discarded at any step.

\begin{lemma} \label{lem:det}
Let $n$ be a positive integer, let $\Lambda$ be a submodule of $A^n$, and let $C_1, \dots, C_k$ be square matrices of size $n$ with entries in $A$ whose rows belong to $\Lambda$.
If the greatest common divisor of the integers $\operatorname{Nm}(\det C_1), \dots, \operatorname{Nm}(\det C_k)$ is $1$, then $\Lambda = A^n$.
\end{lemma}

\begin{proof}
The group $A^n$ is a free $\mathbb{Z}$-module of rank $n(p-1)$.
Let $C$ be one of the matrices $C_i$, and consider the $\mathbb{Z}$-linear map $y \mapsto yC$ from $A^n$ to $A^n$.
After extension of scalars to $\mathbb{C}$, this map is the direct sum of the maps given by the matrices $\tau(C)$, where $\tau$ runs over the $p - 1$ embeddings of $\mathbb{Q}(\omega)$ into $\mathbb{C}$.
Hence its determinant is $\prod_\tau \det \tau(C) = \operatorname{Nm}(\det C)$.
If this integer is not zero, then the image of the map has index $|\operatorname{Nm}(\det C)|$ in $A^n$.
The image is contained in $\Lambda$, so that the index of $\Lambda$ in $A^n$ is finite and divides $\operatorname{Nm}(\det C)$.
By assumption, at least one of the integers $\operatorname{Nm}(\det C_i)$ is not zero, and the index of $\Lambda$ divides each of those that are not zero.
Hence the index is $1$.
\end{proof}

The integer $\operatorname{Nm}(\det C)$ is computed from its residues modulo prime numbers.
For this purpose we need an upper bound.
For $a \in A$, write the image of $a$ in $\mathbb{Z}[\omega]$ as $\sum_{j=0}^{p-1} a_j \omega^j$ with integers $a_j$, and put $t(a) = p \sum_j a_j^2 - (\sum_j a_j)^2$.
This integer is the trace of $|a|^2$ with respect to $\mathbb{Q}(\omega)/\mathbb{Q}$, so that it does not depend on the choice of the integers $a_j$.

\begin{lemma} \label{lem:hadamard}
Let $C = (c_{ij})$ be a square matrix of size $n$ with entries in $A$, and put $t_i = \sum_{j} t(c_{ij})$ for $1 \le i \le n$.
Then
\[
\operatorname{Nm}(\det C)^2 \le \prod_{i=1}^{n} \Bigl( \frac{t_i}{p-1} \Bigr)^{p-1} .
\]
\end{lemma}

\begin{proof}
Let $\tau$ run over the embeddings of $\mathbb{Q}(\omega)$ into $\mathbb{C}$.
By the inequality of Hadamard, $|\det \tau(C)|^2 \le \prod_i \sum_j |\tau(c_{ij})|^2$.
We take the product over $\tau$ and apply the inequality between the geometric mean and the arithmetic mean to each $i$.
Since $\sum_\tau \sum_j |\tau(c_{ij})|^2 = t_i$, the assertion follows.
\end{proof}

\subsection{The lists of elements}

Let $q$ be a split prime.
We describe the prime ideal $\mathfrak{q}_q$ and the test for the condition $\alpha \in s^j(\mathfrak{q}_q)$.
Let $f$ denote the order of $q$ in $(\mathbb{Z}/M\mathbb{Z})^\times$; it divides $p - 1$.
Let $m \in \mathbb{F}_q[t]$ be an irreducible monic polynomial of degree $f$, and let $z$ be an element of order $M$ of the field $\mathbb{F}_q[t]/(m)$.
There is a ring homomorphism $\mathbb{Z}[\zeta] \to \mathbb{F}_q[t]/(m)$ that sends $\zeta$ to $z$.
Its kernel is a prime ideal of $\mathbb{Z}[\zeta]$ above $q$, and we let $\mathfrak{q}_q$ denote the intersection of this kernel with $O_K$.
Since $q$ splits completely in $K$, the images
\[
\rho_i = \sum_{a \in H} z^{g^i a} \qquad (i \in \mathbb{Z}/p\mathbb{Z})
\]
of the Gaussian periods lie in $\mathbb{F}_q$.
Let $\alpha = c + \sum_i a_i \eta_i$ with integers $c$ and $a_i$.
Since $s^j(\mathfrak{q}_q) = \sigma^{-j}(\mathfrak{q}_q)$ and $\sigma^j(\alpha) = c + \sum_i a_i \eta_{i+j}$, we have
\begin{equation} \label{eq:test}
\alpha \in s^j(\mathfrak{q}_q) \iff c + \sum_i a_i \rho_{i+j} = 0 \ \text{ in } \mathbb{F}_q .
\end{equation}

\begin{proposition} \label{prop:datagrh}
Let $p \in \{37, 41, 43\}$, and let $g$, $r$, and $w$ be as in \cref{tab:grh}.
There are $r$ elements of $O_K$ of the form $c + \sum_i a_i \eta_i$ with $c \in \{-2, -1, 0, 1, 2\}$ and $a_i \in \{-1, 0, 1\}$, where the number of indices $i$ with $a_i \ne 0$ is at least $3$ and at most $w$, with the following properties.
Let $T$ denote the set of the split primes that divide the norm of at least one of these elements.
Then the $r$ elements are admissible for $T$, and their vectors generate the $A$-module $A^T$.
The set $T$ has the number of elements stated in \cref{tab:grh}.
Consequently, every prime number in $T$ is settled.
\end{proposition}

\begin{table}[ht]
\centering
\begin{tabular}{ccccccc}
\toprule
$p$ & $g$ & $r$ & $w$ & $|T|$ & $|\{q \in T \mid q < 10^6\}|$ & least elements of $T$ \\
\midrule
$37$ & $2$ & $6039$ & $7$ & $3832$ & $2049$ & $691$, $787$, $1013$ \\
$41$ & $6$ & $7800$ & $8$ & $4434$ & $1865$ & $313$, $487$, $509$ \\
$43$ & $3$ & $5371$ & $8$ & $3528$ & $1832$ & $19$, $367$, $641$ \\
\bottomrule
\end{tabular}
\caption{The lists of elements of \cref{prop:datagrh}.}
\label{tab:grh}
\end{table}

\begin{proof}
The last assertion follows from \cref{prop:span}.
The other assertions are verified by a program with exact integer arithmetic, as follows.

The norm of each element and its factorization are verified as in \cref{subsec:verification}; every prime factor is $p$ or a split prime.
For each $q \in T$, the program chooses a monic polynomial $m$ of degree $f$ and proves that it is irreducible by the criterion of Rabin \cite{Ra80}: the polynomial $m$ is irreducible if and only if $m$ divides $t^{q^f} - t$ and $m$ is coprime to $t^{q^{f/\ell}} - t$ for every prime number $\ell$ dividing $f$.
It chooses an element $z$ of $\mathbb{F}_q[t]/(m)$ with $z^M = 1$ and $z^p \ne 1$, computes the elements $\rho_i$, and checks that they lie in $\mathbb{F}_q$.
For each element $\alpha$ of the list and each split prime $q$ dividing its norm, the program determines the set $J_q$ by \eqref{eq:test} and checks that the number of its elements is the exponent of $q$ in the norm.
Hence $\alpha$ is admissible for $T$.

To prove that the vectors generate $A^T$, the program applies \cref{lem:pivot} repeatedly, each time with a unit of the form $\pm x^j$, and it discards some of the vectors.
After $3712$ steps for $p = 37$, after $4314$ steps for $p = 41$, and after $3408$ steps for $p = 43$, there remain $120$ elements of $T$ and $883$ vectors in $A^{120}$.
Then the program forms square matrices of size $120$ from disjoint sets of these vectors and applies \cref{lem:det}.
For $p = 37$, it uses three matrices: the integers $\operatorname{Nm}(\det C_i)$ are $0$, a number with $2916$ digits, and a number with $3168$ digits.
For $p = 41$, it uses four matrices, and the integers have $3094$, $3591$, $3821$, and $3959$ digits; the greatest common divisor of the first two, and of the first three, is $83$.
For $p = 43$, it uses two matrices, and the integers have $3858$ and $4236$ digits.
In each case, the greatest common divisor of all the integers is $1$.
For $p = 41$, the intermediate value $83$ is the least prime number $\ell$ with $\ell \equiv 1 \pmod{41}$; it is the norm of a prime ideal of $A$, and the first three determinants are divisible by prime ideals of $A$ above $83$.

The integer $\operatorname{Nm}(\det C)$ is computed as follows.
Let $l < 2^{62}$ be a prime number with $l \equiv 1 \pmod{p}$, and let $\theta \in \mathbb{F}_l$ be an element of order $p$.
The residue of $\operatorname{Nm}(\det C)$ modulo $l$ is $\prod_{k=1}^{p-1} \det C(\theta^k)$, where $C(\theta^k)$ denotes the matrix over $\mathbb{F}_l$ obtained from $C$ by the substitution $x \mapsto \theta^k$.
The program uses such prime numbers $l$ until their product exceeds twice the bound of \cref{lem:hadamard}.
\end{proof}

\subsection{The upper bound}

The following is \cite[Lemma 5.2]{Mi14}.

\begin{proposition}[Miller] \label{prop:millergrh}
Let $L$ be a totally real number field of degree $N$, and let $c$ be a positive real number.
Let $S$ be a set of prime numbers each of which splits completely in $L$ into principal prime ideals.
Put
\[
G(c) = \int_0^\infty \frac{1 - e^{-(x/c)^2}}{2} \left( \frac{1}{\sinh(x/2)} + \frac{1}{\cosh(x/2)} \right) dx
\]
and
\[
B' = \frac{\pi}{2} + \gamma + \log 8\pi - \log \operatorname{rd}(L) - G(c) + 2 \sum_{q \in S} \sum_{m=1}^{\infty} \frac{\log q}{q^{m/2}} \, e^{-(m \log q/c)^2} .
\]
Suppose that every nontrivial zero of the Dedekind zeta function of the Hilbert class field of $L$ has real part $1/2$.
If $B' > 0$, then the class number $h$ of $L$ satisfies
\[
h < \frac{2c\sqrt{\pi} \, e^{(c/4)^2}}{N B'} .
\]
\end{proposition}

By \cref{lem:field}, the discriminant of $k_p$ is $p^{2(p-1)}$, so that
\begin{align*}
\log \operatorname{rd}(k_{37}) &= \tfrac{72}{37} \log 37 = 7.0266510732\ldots, \\
\log \operatorname{rd}(k_{41}) &= \tfrac{80}{41} \log 41 = 7.2459942764\ldots, \\
\log \operatorname{rd}(k_{43}) &= \tfrac{84}{43} \log 43 = 7.3474606911\ldots .
\end{align*}

\begin{proposition} \label{prop:boundgrh}
Let $p \in \{37, 41, 43\}$, and let $T$ be the set of \cref{prop:datagrh}.
Then $G(7) \le 0.4709674869$, and for $L = k_p$, $S = T$, and $c = 7$, the double sum $\Sigma'$ and the number $B'$ of \cref{prop:millergrh} satisfy the inequalities of \cref{tab:boundgrh}.
Consequently, if every nontrivial zero of the Dedekind zeta function of the Hilbert class field of $k_p$ has real part $1/2$, then $h(k_p)$ is less than the number in the last column of \cref{tab:boundgrh}.
\end{proposition}

\begin{table}[ht]
\centering
\begin{tabular}{cccc}
\toprule
$p$ & $\Sigma'$ & $B'$ & bound for $h(k_p)$ \\
\midrule
$37$ & $\ge 8.0800417860$ & $\ge 5.9546066$ & $2.4082$ \\
$41$ & $\ge 7.7893310695$ & $\ge 5.4445527$ & $2.3768$ \\
$43$ & $\ge 7.9288497530$ & $\ge 5.4826049$ & $2.2505$ \\
\bottomrule
\end{tabular}
\caption{The bounds of \cref{prop:boundgrh}; here $\Sigma'$ denotes the double sum, including the factor $2$.}
\label{tab:boundgrh}
\end{table}

\begin{proof}
The numbers are computed with the ball arithmetic of Arb \cite{Jo17, Jo18}, as in \cref{sec:bound}.
For the integral, we use the following estimates.
On the interval $[0, \varepsilon]$ with $\varepsilon = 10^{-3}$, the integrand is at most $(x^2/(2c^2))(2/x + 1)$, so that the integral over this interval is at most $(\varepsilon^2/2 + \varepsilon^3/6)/c^2$.
On the interval $[R, \infty)$ with $R = 200$, the integrand is at most $e^{-x/2}(1/(1 - e^{-R}) + 1)$, so that the integral over this interval is at most $2e^{-R/2}(1/(1 - e^{-R}) + 1)$.
The integral over $[\varepsilon, R]$ is enclosed by the rigorous integration of Arb.
The double sum is replaced by the partial sum over all $q \in T$ and $m \le 40$; all terms are positive.
\end{proof}

\begin{corollary} \label{cor:thresholdgrh}
If $h(k_{37}) > 1$, then $h(k_{37}) \ge 149$.
If $h(k_{41}) > 1$, then $h(k_{41}) \ge 83$.
If $h(k_{43}) > 1$, then $h(k_{43}) \ge 173$.
\end{corollary}

\begin{proof}
We argue as in the proof of \cref{cor:threshold}.
Let $\ell$ be a prime number dividing $h(k_p)$, and let $f$ denote the order of $\ell$ modulo $p$; then $\ell \ne p$ and $\ell^f$ divides $h(k_p)$.
The least prime numbers $\ell$ with $\ell \equiv 1 \pmod{p}$ are $149$ for $p = 37$, $83$ for $p = 41$, and $173$ for $p = 43$.
Suppose that $f \ge 2$ and that $\ell^f$ is less than the stated number.
Then $\ell \le 11$ for $p = 37$, $\ell \le 7$ for $p = 41$, and $\ell \le 13$ for $p = 43$.
For $p = 37$, the orders of $2$, $3$, $5$, $7$, and $11$ modulo $37$ are $36$, $18$, $36$, $9$, and $6$, so that $\ell^f > 149$, a contradiction.
For $p = 41$, the orders of $2$, $3$, $5$, and $7$ modulo $41$ are $20$, $8$, $20$, and $40$, so that $\ell^f > 83$, a contradiction.
For $p = 43$, the orders of $2$, $3$, $5$, $7$, $11$, and $13$ modulo $43$ are $14$, $42$, $42$, $6$, $7$, and $21$, so that $\ell^f > 173$, a contradiction.
\end{proof}

\begin{proof}[Proof of \cref{thm:grh}]
By \cref{prop:boundgrh}, we have $h(k_p) < 3$ under the stated hypothesis.
By \cref{cor:thresholdgrh}, it follows that $h(k_p) = 1$.
\end{proof}

\begin{remark} \label{rem:standard}
The argument of this section is a special case of a standard step in the algorithms that compute class groups from relations; see Hafner and McCurley \cite{HM89} and Buchmann \cite{Bu90}.
In these algorithms, a determinant formed from the relations is a multiple of the order of the subgroup of the class group generated by the prime ideals in the relations, and the unit group and the regulator are needed in order to decide whether this multiple is the order itself.
Here the multiple is $1$, so that neither units nor generators of the prime ideals are needed.
The fact that the class group is generated by few prime ideals is replaced by the inequality of \cref{prop:millergrh}.
Since the matrices have entries in $A$ and not in $\mathbb{Z}$, their size is smaller by the factor $p - 1$.
\end{remark}

\begin{remark} \label{rem:entries}
The statement of \cref{prop:datagrh} on the set $T$ does not depend on any hypothesis.
In contrast to \cref{sec:relations}, its proof uses not only the factorizations of the norms but also the sets $J_q$, that is, the information which prime ideals above $q$ contain a given element.
An error in one of these sets may lead to a wrong conclusion.
For this reason, all sets $J_q$ of the three lists were computed a second and a third time with independent programs: once with ring homomorphisms $O_K \to \mathbb{F}_q$ that are constructed without the field $\mathbb{Q}(\zeta)$, from the multiplication table of the basis $1, \eta_0, \dots, \eta_{p-2}$ of $O_K$, and once with the ideal arithmetic of PARI/GP \cite{PARI}.
The results agree.
The generation of $A^T$ was also verified with a second program that uses other pivots and other matrices in \cref{lem:det}.
\end{remark}

\begin{remark} \label{rem:control}
We ran the same program on cyclic fields of prime degree $d$ and prime conductor $\ell$ whose class groups are known from PARI/GP \cite{PARI}; in this case $A = \mathbb{Z}[x]/(1 + \dots + x^{d-1})$.
For $(\ell, d) = (13, 3)$, $(577, 3)$, $(401, 5)$, and $(71, 7)$, the class number is $1$, and the greatest common divisor of \cref{lem:det} was $1$.
For $(\ell, d) = (163, 3)$, $(349, 3)$, $(313, 3)$, $(941, 5)$, and $(491, 7)$, the class numbers are $4$, $4$, $7$, $16$, and $8$, and in our runs the greatest common divisor was equal to the class number.
\end{remark}

\begin{remark} \label{rem:searchgrh}
The elements of \cref{prop:datagrh} were found as follows; this is not part of the proof.
We computed the norms of the elements $c + \sum_i a_i \eta_i$ with $c \in \{-2, \dots, 2\}$ and $a_i \in \{-1, 0, 1\}$ with a given number of nonzero coefficients $a_i$, one element for each orbit under the Galois group and the sign.
We kept the elements whose norm is a product of a power of $p$, of split primes below $10^6$, and of at most one further prime number below $10^9$, and then we removed the elements with a prime factor that occurs in no other element.
For $p = 37$, we examined about $1.0 \cdot 10^8$ elements with at most $7$ nonzero coefficients; this took about $6$ minutes on two cores.
For $p = 41$ and $p = 43$, we examined about $1.05 \cdot 10^9$ and $7.4 \cdot 10^8$ elements with at most $8$ nonzero coefficients.
The verification takes between $10$ and $35$ minutes for each field on one core.
For $p = 37$, the function \texttt{bnfinit} of PARI/GP did not finish within $355$ minutes in our attempt.
\end{remark}

\section{Remarks} \label{sec:remarks}

\begin{remark} \label{rem:first}
For $k_{7^2}$, we also carried out the method of \cite{Mi15b} directly.
For each of the $28381$ split primes $q \le 2.2 \cdot 10^7$, an element $\alpha_q \in O_K$ with $|N_{K/\mathbb{Q}}(\alpha_q)| = q$ was computed with the functions \texttt{bnfinit} and \texttt{bnfisprincipal} of PARI/GP and then verified independently.
With these prime numbers alone, \cref{prop:miller} gives $h(k_{7^2}) < 160.7121$ at $c = 80$.
This computation took about $100$ hours on one core, and the elements $\alpha_q$ occupy $187$ megabytes; the list of relations of \cref{prop:data7} occupies $2.4$ megabytes.
\end{remark}

\begin{remark} \label{rem:19}
In $k_{7^2}$, the prime number $19$ alone contributes more than $0.6222$ to $\Sigma(S_7, 80)$, whereas the next split prime $1373$ contributes less than $0.0209$.
In $k_{3^4}$, the least split prime is $487$.
This is the reason why the set $S_3$ has to be much larger than the set $S_7$, although the root discriminant of $k_{3^4}$ is smaller.
\end{remark}

\begin{remark} \label{rem:ramified}
\Cref{prop:miller} is derived in \cite{Mi14} from the identity \cite[(4.1)]{Mi14}, which contains a sum over all prime ideals of the Hilbert class field of $L$, and every term of this sum is positive.
The terms that belong to other principal prime ideals of $L$ may therefore be kept as well; this is done for a ramified prime ideal in \cite[Section 6]{Mi14}.
For $L = k_{p^n}$, the prime ideal above $p$ is principal of norm $p$ by \cref{lem:field}, so that it splits completely in the Hilbert class field.
Keeping its terms adds
\[
\frac{2}{N} \sum_{m=1}^{\infty} \frac{\log p}{p^{m/2}} F(m \log p)
\]
to the number $B$.
With this term, the bounds of \cref{prop:bound} improve to $h(k_{3^4}) < 65.9166$ at $c = 80$, $h(k_{7^2}) < 93.9394$ at $c = 80$, and $h(k_{29}) < 39.1488$ at $c = 60$.
The proof of \cref{thm:main} does not use this refinement.
\end{remark}

\begin{remark} \label{rem:fkm}
Fukuda, Komatsu, and Morisawa \cite{FKM14} proved that no prime number less than $10^9$ divides $h(k_{3^n})$ for any $n$; see \cite[Theorem 5]{Mi15b}.
The proof of \cref{thm:main} does not use this result.
\end{remark}

\begin{remark} \label{rem:23}
The equality $h(k_{23}) = 1$ is claimed in \cite{Ch26}.
There, relations for $99167$ of the split primes of $k_{23}$ below $5 \cdot 10^7$ and for $55$ larger split primes are obtained from an exhaustive search over the elements with coefficients in $\{-1, 0, 1\}$, and \cref{prop:miller} with $c = 50$ gives $h(k_{23}) < 37.72$; according to \cite{Ch26}, the search took about $25$ core-hours.
We confirmed this equality with the method of the present note.
For $p^n = 23$, we have $M = 529$, the group $H$ has order $22$, and we take $g = 5$.
There are $130425$ split primes of $k_{23}$ below $5 \cdot 10^7$, and all of them are settled: for the $150$ split primes below $3 \cdot 10^4$, we use elements whose norm has absolute value $q$, computed as in \cref{rem:41}, and for the others we use \cref{prop:range} with $A = 3 \cdot 10^4$.
With $c = 60$, \cref{prop:miller} gives $B \ge 0.3068684$ and $h(k_{23}) < 30.1354$, The argument of \cref{cor:threshold} shows that $h(k_{23})$ is either $1$ or at least $47$: the least prime number $\ell \equiv 1 \pmod{23}$ is $47$, and the orders of $2$, $3$, and $5$ modulo $23$ are $11$, $11$, and $22$.
Our search took about $4$ minutes on two cores.
\end{remark}

\begin{remark} \label{rem:limit}
In \cref{prop:miller}, the inequality $B > 0$ requires that the sum over $S$ exceeds $\log \operatorname{rd}(L) + I(c) - \pi/2 - \gamma - \log 8\pi$, and the prime numbers $q \le X$ that split completely in $L$ contribute about $(4/N) \log X$ to this sum.
Hence the number of prime numbers that are needed grows quickly with the degree and with the root discriminant.
By a numerical estimate, the field $k_{31}$ requires the split primes up to at least $10^{11}$.
Under the generalized Riemann hypothesis, \cref{prop:millergrh} requires only a small set $S$; see \cref{sec:grh}.
\end{remark}

\begin{thebibliography}{99}

\bibitem{Ba69} H.~Bauer, Numerische Bestimmung von Klassenzahlen reeller zyklischer Zahlk\"orper, J. Number Theory \textbf{1} (1969), 161--162.

\bibitem{Bu90} J.~Buchmann, A subexponential algorithm for the determination of class groups and regulators of algebraic number fields, in: S\'eminaire de Th\'eorie des Nombres, Paris 1988--1989, Progr. Math. \textbf{91}, Birkh\"auser Boston, Boston, MA, 1990, pp.~27--41.

\bibitem{Ch26} J.~Chakkour, \texttt{b23-class-number}, software repository, 2026, \url{https://github.com/johnchakkour/b23-class-number}.

\bibitem{Co12} J.~Coates, The enigmatic Tate--Shafarevich group, in: Fifth International Congress of Chinese Mathematicians, AMS/IP Stud. Adv. Math. \textbf{51}, Amer. Math. Soc., Providence, RI, 2012, pp.~43--50.

\bibitem{Co60} H.~Cohn, A numerical study of Weber's real class number of calculation. I, Numer. Math. \textbf{2} (1960), 347--362.

\bibitem{FK11} T.~Fukuda and K.~Komatsu, Weber's class number problem in the cyclotomic $\mathbb{Z}_2$-extension of $\mathbb{Q}$, III, Int. J. Number Theory \textbf{7} (2011), 1627--1635.

\bibitem{FKM14} T.~Fukuda, K.~Komatsu, and T.~Morisawa, Weber's class number one problem, in: Iwasawa Theory 2012, Contrib. Math. Comput. Sci. \textbf{7}, Springer, Heidelberg, 2014, pp.~221--226.

\bibitem{FLINT} The FLINT team, FLINT: Fast Library for Number Theory, \url{https://flintlib.org}.

\bibitem{HM89} J.~L.~Hafner and K.~S.~McCurley, A rigorous subexponential algorithm for computation of class groups, J. Amer. Math. Soc. \textbf{2} (1989), 837--850.

\bibitem{HH09} K.~Horie and M.~Horie, The ideal class group of the $\mathbb{Z}_p$-extension over the rationals, Tohoku Math. J. \textbf{61} (2009), 551--570.

\bibitem{Iw56} K.~Iwasawa, A note on class numbers of algebraic number fields, Abh. Math. Sem. Univ. Hamburg \textbf{20} (1956), 257--258.

\bibitem{Jo17} F.~Johansson, Arb: efficient arbitrary-precision midpoint-radius interval arithmetic, IEEE Trans. Comput. \textbf{66} (2017), 1281--1292.

\bibitem{Jo18} F.~Johansson, Numerical integration in arbitrary-precision ball arithmetic, in: Mathematical Software -- ICMS 2018, Lecture Notes in Comput. Sci. \textbf{10931}, Springer, Cham, 2018, pp.~255--263.

\bibitem{Ma78} J.~M.~Masley, Class numbers of real cyclic number fields with small conductor, Compositio Math. \textbf{37} (1978), 297--319.

\bibitem{Mi14} J.~C.~Miller, Class numbers of totally real fields and applications to the Weber class number problem, Acta Arith. \textbf{164} (2014), 381--398.

\bibitem{Mi14b} J.~C.~Miller, Class numbers of real cyclotomic fields of composite conductor, LMS J. Comput. Math. \textbf{17} (2014), suppl. A, 404--417.

\bibitem{Mi15b} J.~C.~Miller, Class numbers in cyclotomic $\mathbb{Z}_p$-extensions, J. Number Theory \textbf{150} (2015), 47--73.

\bibitem{MO13} T.~Morisawa and R.~Okazaki, Mahler measure and Weber's class number problem in the cyclotomic $\mathbb{Z}_p$-extension of $\mathbb{Q}$ for odd prime number $p$, Tohoku Math. J. \textbf{65} (2013), 253--272.

\bibitem{PARI} The PARI~Group, PARI/GP, Univ. Bordeaux, \url{https://pari.math.u-bordeaux.fr}.

\bibitem{Ra80} M.~O.~Rabin, Probabilistic algorithms in finite fields, SIAM J. Comput. \textbf{9} (1980), 273--280.

\bibitem{vdL82} F.~J.~van der Linden, Class number computations of real abelian number fields, Math. Comp. \textbf{39} (1982), 693--707.

\bibitem{Wa97} L.~C.~Washington, Introduction to Cyclotomic Fields, 2nd ed., Grad. Texts in Math. \textbf{83}, Springer, New York, 1997.

\bibitem{We86} H.~Weber, Theorie der Abel'schen Zahlk\"orper, Acta Math. \textbf{8} (1886), 193--263.

\end{thebibliography}

\end{document}
